\section{Validierung}
\label{sec:validierung}

% VERANTWORTLICH: Linus
% Aufgabenpunkt 4.
% ROHFASSUNG. Abschnitt 5.1 steht, 5.2 und 5.3 warten auf die
% Gruppenentscheidung zur Datengrundlage (siehe doku/annahmen.md,
% offener Punkt 1).

\subsection{Vorgehen und Fehlermasse}

Die Vorlesung nennt drei Kriterien fuer die Validierung eines Modells:
Reproduzierbarkeit, Praediktion und innere Struktur. Auf das vorliegende
Modell angewandt bedeutet das:

\begin{itemize}
  \item \emph{Reproduzierbarkeit.} Bei gleichen Eingangsdaten liefert die
        Implementierung dasselbe Ergebnis. Da alle Parameter an einer
        einzigen Stelle liegen und der Loeser deterministisch arbeitet, ist
        dieses Kriterium konstruktiv erfuellt. Der Nachweis der
        Loesungsunabhaengigkeit in Abschnitt~\ref{sec:solver} gehoert
        ebenfalls hierher.
  \item \emph{Praediktion.} Das Modell muss Messwerte vorhersagen, die
        nicht in seine Aufstellung eingeflossen sind. Dies ist der Kern
        dieses Kapitels.
  \item \emph{Innere Struktur.} Nicht nur das Ergebnis, auch die einzelnen
        Terme muessen physikalisch sinnvoll sein. Zwei Beobachtungen stuetzen
        das: In den Nachtstunden liegt die Modultemperatur um zwei bis drei
        Kelvin \emph{unter} der Lufttemperatur. Das ist kein Rechenfehler,
        sondern die Abstrahlung gegen den nach Swinbank deutlich kaelteren
        Himmel, und es bestaetigt den Strahlungsterm unabhaengig vom
        Gesamtergebnis. Ausserdem schliesst die Energiebilanz ueber eine
        Woche auf ein Hundertstel Prozent, siehe
        Abschnitt~\ref{sec:solver}.
\end{itemize}

\paragraph{Wahl der Fehlermasse.}
Verwendet werden der mittlere absolute Fehler (MAE), die Wurzel des
mittleren quadratischen Fehlers (RMSE) und der mittlere systematische Fehler
(MBE), jeweils in Kelvin. Die drei Masse beantworten verschiedene Fragen:
der MAE gibt die typische Abweichung an, der RMSE gewichtet einzelne grosse
Ausreisser staerker, und der MBE zeigt als einziger die \emph{Richtung} der
Abweichung. Ein Modell kann einen kleinen RMSE aufweisen und trotzdem
systematisch zu tief liegen; erst der MBE macht das sichtbar. Das Residuum
ist als $r = T_\mathrm{sim} - T_\mathrm{mess}$ definiert, ein negativer MBE
bedeutet also eine Unterschaetzung. Diese Konvention entspricht der von
Tuncel et al.~\cite{tuncel2020}, deren Werte damit unmittelbar
vergleichbar sind.

\paragraph{Getrennte Auswertung fuer Tag und Nacht.}
Die Fehlermasse werden zusaetzlich getrennt fuer die Tages- und die
Nachtstunden angegeben. Der Grund dafuer laesst sich am Referenzpaper selbst
ablesen: Tuncel et al. berichten fuer denselben Datensatz einen MAE von
$0{,}90\,\mathrm{K}$ ueber den gesamten Zeitraum, aber $2{,}61\,\mathrm{K}$
allein fuer die Tagesstunden. Nachts liegt das Modul nahe an der
Umgebungstemperatur, die Dynamik ist gering und die Abweichung klein. Wer
ueber volle Tage mittelt, verduennt die Abweichung der Mittagsstunden mit
vielen unauffaelligen Nachtwerten und erhaelt eine Zahl, die guenstiger
aussieht, als das Modell ist. Als Grenze gilt eine Globalstrahlung von
$20\,\mathrm{W/m^2}$; das Referenzpaper definiert die Grenze nicht, die
Festlegung ist daher in der Annahmentabelle dokumentiert.

\paragraph{Kriterium fuer eine bestandene Validierung.}
% TODO In der Gruppe bestaetigen, siehe doku/annahmen.md, offener Punkt 2.
Das Kriterium wurde festgelegt, \emph{bevor} Ergebnisse vorlagen. Wer die
Schwelle erst nach dem Ergebnis waehlt, waehlt die guenstigste. Als Grenze
gilt ein MAE der Tagesstunden von $5\,\mathrm{K}$. Dieser Wert ist aus der
Literatur begruendet: Tuncel et al. erreichen mit dem vollstaendigen
Nusselt-Ansatz und vier unterschiedenen Stroemungsregimen tagsueber einen MAE
von $2{,}61\,\mathrm{K}$. Das hier verwendete Basismodell fasst freie und
erzwungene Konvektion in einem linearen Ansatz zusammen und kann diesen Wert
systematisch nicht unterbieten. Ein Kriterium in der genannten
Groessenordnung ist damit sachlich hergeleitet und nicht nachtraeglich
passend gewaehlt.

\subsection{Datengrundlage}
\label{sec:validierung_daten}

% TODO Dieser Abschnitt haengt an der Gruppenentscheidung, siehe
% doku/annahmen.md, offener Punkt 1. Der folgende Absatz haelt den Befund
% fest, der die urspruengliche Planung hinfaellig gemacht hat.

Als Referenz war Abbildung~1(a) aus Tuncel et al.~\cite{tuncel2020}
vorgesehen, die gemessene und modellierte Modultemperatur ueber fuenf Tage
zeigt. Bei der Vorbereitung der Digitalisierung stellte sich heraus, dass
die Abbildung ausschliesslich die Zielgroesse enthaelt. Globalstrahlung,
Umgebungstemperatur und Windgeschwindigkeit, also genau die drei
Eingangsgroessen des Modells, sind weder dort noch an anderer Stelle der
Arbeit als Zahlenwerte veroeffentlicht; Abschnitt~2.4 der Arbeit haelt nur
fest, dass sie stuendlich an der Wetterstation der Anlage erfasst wurden.
Aus der Abbildung laesst sich damit gewinnen, was das Modell vorhersagen
soll, nicht aber das, womit es rechnet.

\subsection{Ergebnisse}

\begin{figure}[htbp]
  \centering
  % \includegraphics[width=\textwidth]{validierung_temperatur}
  \caption{Simulierte und im Referenzpaper gemessene Modultemperatur ueber
           fuenf Tage. Die Abweichung ist an den Mittagsstunden der klaren
           Tage am groessten.}
  \label{fig:validierung}
\end{figure}

% TODO Zahlen einsetzen, sobald die Datengrundlage geklaert ist.
% Die Konsolenausgabe von plot_validation.m liefert die Werte fuer die
% Bildunterschrift bereits fertig formatiert.

\subsection{Diskussion der Abweichungen}

% TODO Quantifizieren statt entschuldigen: welcher Effekt erklaert wie viel
% Kelvin? Zu behandeln sind mindestens:
%   - fehlende Angaben zu den Eingangsgroessen im Referenzpaper
%   - linearer Konvektionsansatz statt Nusselt-Ansatz
%   - Annahme A_conv = 2A ueber die Einbausituation
%   - ein Temperaturknoten statt Schichtmodell
%   - Ableseungenauigkeit beim Digitalisieren
